\documentclass[11pt,a4paper]{article}

\usepackage[utf8]{inputenc}
\usepackage[T1]{fontenc}
\usepackage{amsmath,amssymb,amsthm}
\usepackage{mathtools}
\usepackage{hyperref}
\usepackage{geometry}
\geometry{margin=1in}
\usepackage{enumitem}

\newtheorem{definition}{Definition}[section]
\newtheorem{proposition}[definition]{Proposition}
\newtheorem{remark}[definition]{Remark}
\newtheorem{example}[definition]{Example}

\DeclareMathOperator{\hash}{hash}
\newcommand{\Cov}{\mathcal{C}}
\newcommand{\N}{\mathbb{N}}

\title{Coverings and Connections:\\ A Mathematical Foundation for \texttt{covmap}}
\author{}
\date{\today}

\begin{document}
\maketitle

\begin{abstract}
\texttt{covmap} links two repositories $A$ and $B$ by partitioning each
into content-addressed \emph{segments} and relating those segments
across repositories. This note gives precise, minimal definitions of
repositories, segments, coverings, sharing of segments across
coverings, and connections, then describes the command interface that
realizes them. The implementation deliberately ships no segmentation
\emph{rule}: every covering is built up by hand, segment by segment,
from a single whole-file bootstrap.
\end{abstract}

\tableofcontents

\section{Segments and coverings}\label{sec:cov}

\subsection{Repositories}

\begin{definition}[Repository]\label{def:repo}
A \emph{repository} $A$ is a finite set of files, each a finite
sequence of lines. We write $|f|_A$ for the number of lines of file
$f$ in $A$, and $\mathrm{content}_A(f, i)$ for the text of line $i$ of
file $f$. \texttt{covmap} considers only files whose bytes are valid
UTF-8; other files are skipped on bootstrap and never appear in any
covering.
\end{definition}

\subsection{Segments}

A segment is a contiguous range of lines inside a single file.

\begin{definition}[Segment]\label{def:segment}
Fix a set of \emph{labels} $L$. A \emph{segment} of $A$ is a
quadruple $s = (f, a, b, \lambda)$ where $f$ is a file of $A$,
$1 \le a \le b \le |f|_A$, and $\lambda \in L$. Its \emph{lines} are
\[
  \mathrm{lines}(s) \;=\; \{(f, i) \mid a \le i \le b\},
\]
its \emph{content} is the concatenation
\[
  \mathrm{content}_A(s)
  \;=\;
  \mathrm{content}_A(f, a) \;\Vert\;
  \mathrm{content}_A(f, a+1) \;\Vert\;
  \cdots \;\Vert\;
  \mathrm{content}_A(f, b),
\]
and its \emph{label} is $\mathrm{label}(s) = \lambda$. We fix a hash
function $h$ and define the \emph{content hash} of $s$ as
\[
  \hash(s) \;=\; \mathrm{trunc}_{16}\bigl(h\bigl(\lambda \;\Vert\; \texttt{NUL} \;\Vert\; \mathrm{content}_A(s)\bigr)\bigr),
\]
the first 16 hex characters (64 bits) of the SHA-256 of the labelled
content. The truncated hash is the segment's identity.
\end{definition}

\begin{remark}[Why 64-bit identity]
With $n$ segments in play, the birthday-bound collision probability
is approximately $n^2 / 2^{65}$. For repositories with up to a few
tens of thousands of segments, this is negligible; the saving in
display and storage compared to full SHA-256 is significant.
\end{remark}

The label is part of identity: two segments with identical content
but different labels have different hashes and are therefore
distinguishable as endpoints of a connection. \texttt{covmap}
bootstraps every segment with the label \texttt{unknown}; typical
re-labels include \emph{ignored}, \emph{reviewed}, and \emph{mapped},
used to record per-segment status while preserving content-based
sharing among segments that agree in both content and label.

Two segments $s_1 = (f, a_1, b_1, \lambda_1)$ and
$s_2 = (f, a_2, b_2, \lambda_2)$ in the same file are
\emph{adjacent} if $b_1 + 1 = a_2$, and they \emph{overlap} if their
line ranges intersect (regardless of label).

\subsection{Coverings}

\begin{definition}[Covering]\label{def:covering}
Let $F \subseteq \mathrm{files}(A)$. A \emph{covering of $A$ on $F$}
is a finite set of segments $\Cov = \{s_1, \dots, s_n\}$ such that,
for every file $f \in F$, the segments of $\Cov$ that lie in $f$ form
a partition of $f$ into contiguous blocks:
\begin{enumerate}[label=(C\arabic*),leftmargin=2.5em]
\item \textbf{Disjointness:} no two segments of $\Cov$ overlap;
\item \textbf{Completeness:} every line $(f, i)$ with $f \in F$ and
      $1 \le i \le |f|_A$ lies in some segment of $\Cov$;
\item \textbf{Confinement:} every segment of $\Cov$ lies in some
      file of $F$.
\end{enumerate}
We call $F = F(\Cov)$ the \emph{covered set} of $\Cov$.
\end{definition}

A covering need not partition the whole repository. It declares which
files it accepts responsibility for. New files appearing in $A$ are
covered by re-running the bootstrap on $\Cov$ (\S\ref{sec:bootstrap}),
which is additive and never disturbs existing segments.

Equivalently, for each file $f \in F$ a covering supplies a sequence
of cut points $1 = c_0 < c_1 < \cdots < c_{k_f} = |f|_A + 1$ together
with a label $\lambda_j \in L$ for each block, and the segments in
$f$ are $(f, c_{j-1}, c_j - 1, \lambda_j)$.

\begin{proposition}[Block lookup]\label{prop:lookup}
For any covering $\Cov$ and any line $(f, i)$ of $A$ with
$f \in F(\Cov)$, exactly one segment of $\Cov$ contains $(f, i)$.
\end{proposition}

\subsection{Multiple coverings and sharing}\label{sec:sharing}

A repository may carry several coverings simultaneously. Different
coverings often agree on some blocks. Because identity is
content-addressed, agreement is automatic: two coverings $\Cov_1$ and
$\Cov_2$ \emph{share} a segment $s$ exactly when $s \in \Cov_1$,
$s \in \Cov_2$, and $\hash(s)$ is the same in both.

\begin{definition}[Segment store]\label{def:store}
The \emph{segment store} for $A$ is a partial map
\[
  \sigma_A : \{\text{hashes}\} \rightharpoonup \{\text{segments of } A\},
\]
populated lazily as coverings are computed. A covering $\Cov$ is
represented as a finite list of $(\hash(s), s)$ entries, and the
segments themselves are fetched from the store by hash.
\end{definition}

\begin{definition}[DAG of coverings and segments]\label{def:dag}
Given a family of coverings $\Cov_1, \dots, \Cov_n$ of $A$, define a
directed graph with one node per covering, one node per distinct
segment hash that appears in some $\Cov_i$, and an edge $\Cov_i \to s$
whenever $s \in \Cov_i$.
\end{definition}

\begin{proposition}\label{prop:dag-acyclic}
This graph is acyclic, and a segment that belongs to several
coverings is represented exactly once.
\end{proposition}

\section{Connections}\label{sec:conn}

\begin{definition}[Connection]\label{def:connection}
Let $\Cov^A$ be a covering of $A$ and $\Cov^B$ a covering of $B$. A
\emph{connection} from $\Cov^A$ to $\Cov^B$ is a relation
\[
  \kappa \;\subseteq\; \Cov^A \times \Cov^B.
\]
\end{definition}

A pair $(s, t) \in \kappa$ records that segment $s$ of $A$ is linked
to segment $t$ of $B$. Connections are stored separately from
coverings and refer to their endpoints by hash.

\begin{definition}[Coverage]\label{def:coverage}
The \emph{source coverage} of $\kappa$ is
\[
  \rho_{\mathrm{src}}(\kappa)
  \;=\;
  \frac{|\{s \in \Cov^A \mid \exists t.\ (s, t) \in \kappa\}|}{|\Cov^A|}
  \;\in\; [0, 1].
\]
The \emph{target coverage} is the symmetric quantity for $\Cov^B$. A
connection is \emph{source-total} if $\rho_{\mathrm{src}}(\kappa) = 1$.
\end{definition}

\begin{definition}[Operations]\label{def:conn-ops}
Connections support
\begin{itemize}[leftmargin=2em]
\item \emph{inverse:} $\kappa^{-1} = \{(t, s) \mid (s, t) \in \kappa\}$;
\item \emph{composition:} for
      $\kappa \subseteq \Cov^A \times \Cov^B$ and
      $\lambda \subseteq \Cov^B \times \Cov^C$,
      $\lambda \circ \kappa
       = \{(s, u) \mid \exists t.\ (s, t) \in \kappa,\ (t, u) \in \lambda\}$;
\item \emph{union} of two connections with the same endpoints.
\end{itemize}
\end{definition}

\section{Complexity of a covering}\label{sec:complexity}

A covering can be characterized by several quantitative measures.

Let $\ell(s) = b - a + 1$ be the \emph{length} of a segment
$s = (f, a, b)$, and let $N_\Cov = \sum_{s \in \Cov} \ell(s)$ be the
total number of covered lines.

\subsection{Cardinality and granularity}

\begin{definition}[Cardinality, granularity]\label{def:size}
The \emph{cardinality} of $\Cov$ is $|\Cov|$, and its
\emph{granularity} is
\[
  g(\Cov) \;=\; \frac{|\Cov|}{N_\Cov} \;\in\; (0, 1].
\]
\end{definition}

A bootstrap covering (one segment per covered file) has $g$ close to
$0$; the one-line-per-segment covering has $g = 1$.

\subsection{Partition entropy}

\begin{definition}[Partition entropy]\label{def:entropy}
Let $p_\Cov(s) = \ell(s)/N_\Cov$. The \emph{entropy} of $\Cov$ is
\[
  H(\Cov) \;=\; -\sum_{s \in \Cov} p_\Cov(s) \log_2 p_\Cov(s).
\]
\end{definition}

Entropy ranges from $0$ (a single block covering one one-line file)
to $\log_2 N_\Cov$ (one block per line). Coverings with equal-length
blocks maximize entropy at fixed cardinality.

\subsection{Edit stability and drift}

A covering exists to be reused after the repository is edited. We
measure how stable a covering is under edits.

\begin{definition}[File-level drift]\label{def:drift-file}
Let $\Cov$ be a covering of $A$ and let $A'$ be a repository obtained
from $A$ by an edit. A file $f \in F(\Cov)$ is \emph{drifted} under
this edit when at least one segment $s = (f, a, b, \lambda) \in \Cov$
satisfies $\mathrm{content}_{A'}(s) \neq \mathrm{content}_A(s)$.
\end{definition}

\begin{definition}[Boundary drift]\label{def:drift-bound}
A drifted file $f$ has \emph{boundary drift} if some segment of
$\Cov$ in $f$ has $b > |f|_{A'}$, i.e., the recorded line range no
longer fits the file.
\end{definition}

The two flavours of drift correspond to two repair operations
(\S\ref{sec:interface}): \texttt{recut} for file-level drift with
intact boundaries, \texttt{merge} + \texttt{split} for boundary
drift.

\begin{definition}[Broken pair]\label{def:broken}
A pair $(s, t) \in \kappa$ is \emph{broken} relative to current
$\Cov^A$ and $\Cov^B$ when $\hash(s) \notin \hash(\Cov^A)$ or
$\hash(t) \notin \hash(\Cov^B)$.
\end{definition}

A pair becomes broken precisely when its endpoint is rewritten
(\texttt{recut}, \texttt{label}, \texttt{split}, \texttt{merge}) or
removed.

\section{An interface for \texttt{covmap}}\label{sec:interface}

The model above suggests a small command surface, organized
\texttt{git}-style around an object store, lightweight named refs,
and human-readable text refs.

\subsection{State on disk}

A repository under \texttt{covmap} acquires a \texttt{.covmap/}
directory:
\begin{verbatim}
.covmap/
  HEAD                default covering name
  ignore              optional gitignore-lite patterns
  objects/<hh>/<rest> segments, addressed by hash
  coverings/<name>    one file per covering (covered files + segments)
  connections/<name>  one file per connection (metadata + pairs)
\end{verbatim}

\subsection{Addressing}

Every object has a stable address. Three address forms appear
throughout the interface and can be mixed freely.

\begin{tabbing}
\hspace{4.2cm}\=\kill
\texttt{<hash>}            \> a segment, by full or short hex prefix\\
\texttt{$C$:$h$}           \> a segment $h$ in covering $C$\\
\texttt{$K$}               \> a connection by name
\end{tabbing}

A bare \texttt{$h$} is short for \texttt{HEAD:$h$}. Where a covering
is expected, the literal \texttt{HEAD} resolves to the default
covering recorded in \texttt{.covmap/HEAD}. Hash prefixes follow
\texttt{git}-style disambiguation: any unique prefix works,
ambiguity is an error.

\subsection{Bootstrap and refinement}\label{sec:bootstrap}

There is no segmentation rule. The interface gives the user one
bootstrap operation and two refinement operations.

\begin{itemize}[leftmargin=2em,itemsep=2pt]
\item \texttt{covmap init}\\
      Create \texttt{.covmap/} in the current directory.
\item \texttt{covmap cover [$C$]}\\
      For every readable UTF-8 file under the repo not already
      covered by $C$, add a single segment $(f, 1, |f|_A,
      \texttt{unknown})$ to $C$. Idempotent: re-running \texttt{cover}
      adds segments only for new files. Existing segments are never
      touched, so any user-driven refinement is preserved.
\item \texttt{covmap split $C$:$h$ $c$ [$\lambda_t$] [$\lambda_b$]}\\
      Replace the segment with hash $h$ by two adjacent segments
      $(f, a, c-1, \lambda_t)$ and $(f, c, b, \lambda_b)$, where
      $a \le c \le b$ and $a < c$. Both new labels default to
      \texttt{unknown}.
\item \texttt{covmap merge $C$:$h_1$ $C$:$h_2$ [$\lambda$]}\\
      Replace two adjacent segments by their union, labelled
      $\lambda$ (default \texttt{unknown}).
\item \texttt{covmap recut $C$:$h$}\\
      Refresh the segment's hash from current disk bytes. Keeps
      $(f, a, b, \lambda)$; rewrites the content from the file as it
      is now. Errors if the segment no longer fits the file (boundary
      drift); use \texttt{merge}+\texttt{split} in that case.
\item \texttt{covmap label $C$:$h$ $\lambda$}\\
      Replace the label of one segment. Because labels are part of
      identity, this rewrites the entry in $C$; old hashes remain in
      the object store so any connection that referenced them keeps
      resolving.
\end{itemize}

These five operations are the only way to change a covering after
\texttt{cover}. Together they are complete: any covering reachable
from the bootstrap by repeated splits / merges / recuts / relabels
is the most general covering one can ask for.

\subsection{Inspection, connections, repair}

\begin{itemize}[leftmargin=2em,itemsep=2pt]
\item \texttt{covmap status [$C$]}\\
      One-line summary per covering: cardinality, line count,
      entropy. A drift line reports the number of files whose bytes
      have changed since their segments were recorded. Broken pair
      counts are reported per connection.
\item \texttt{covmap show <addr>}\\
      Print the segment, covering, or connection at \texttt{<addr>}.
\item \texttt{covmap ls [-l] $C$}\\
      List segments of $C$ in file-and-line order.
\item \texttt{covmap grep [-l] [-c $C$] <pattern>}\\
      Print segments whose label or content matches \texttt{<pattern>}.
\item \texttt{covmap diff $C_1$ $C_2$}\\
      Segments only in $C_1$, only in $C_2$, and shared.
\item \texttt{covmap connect $C_A$ $\langle\text{repo}_B\rangle$:$C_B$ [$K$]}\\
      Create an empty connection between two coverings.
\item \texttt{covmap link $K$ $h_A$ $h_B$ [$\nu$]}\\
      \texttt{covmap unlink $K$ $h_A$ $h_B$}\\
      Add or remove a single pair, optionally annotated with note
      $\nu$.
\item \texttt{covmap repair $K$}\\
      Drop every pair of $K$ whose endpoint is no longer in its
      covering.
\item \texttt{covmap fsck}\\
      Verify that every covering is a partition of its covered set
      and that every connection's endpoints resolve.
\end{itemize}

\subsection{Plumbing}

Lower-level commands meant for scripts.

\begin{itemize}[leftmargin=2em,itemsep=2pt]
\item \texttt{covmap hash-segment <file> <a> <b> [<label>] [-w]}
\item \texttt{covmap cat-segment <hash>}
\item \texttt{covmap ls-segments $C$}
\item \texttt{covmap ls-pairs $K$}
\item \texttt{covmap rev-parse <addr>}
\item \texttt{covmap fsck --broken-pairs}
\end{itemize}

\subsection{Design discipline}

\begin{itemize}[leftmargin=2em]
\item Coverings and connections are independent artifacts; a
      connection refers to its coverings only by hash.
\item Every state-changing operation is explicit. \texttt{cover} is
      additive only. Refining a covering requires \texttt{split},
      \texttt{merge}, \texttt{recut}, or \texttt{label} —
      \texttt{covmap} never silently re-segments existing files.
\item The plumbing/porcelain split is hard: every porcelain command
      is implementable as a shell pipeline of plumbing commands.
\item Hashes are the lingua franca. Every command accepts them on
      input and prints them on output, so commands compose without
      glue code.
\end{itemize}

\end{document}
